\section{The Security Proof}\label{sec:proof}

\begin{theorem}[\m{model_sliced_NI}]\label{thm:main}
If the scheduler and the public user process are non-interfering at
\m{public_equiv} on input and output, and the private user process is
non-interfering at \m{private_equiv} on input and output, then
\begin{multline*}
\mathsf{model\_sliced}\;\mathit{runtime}\;\mathit{runs}\;p_{pub}\;p_{priv}\;p_{sched}\\
  \in \NI(\mathsf{in\_equiv},\, \mathsf{out\_equiv}).
\end{multline*}
\end{theorem}

\begin{corollary}[\m{model_sliced_userview_NI}]
Under the same hypotheses, \m{model_sliced_userview} is non-interfering at
\m{out_equiv_userview}.
\end{corollary}

The corollary follows from \cref{thm:main} by \m{map_NI}, since \m{parse_output} maps
\m{out_equiv}-related outputs to \m{out_equiv_userview}-related ones. The rest of this
section is about \cref{thm:main}.

\subsection{The equivalences}\label{sec:equivs}

\paragraph{Input} \m{in_equiv} makes the disk interrupt, and only the disk interrupt,
unobservable at $\Lo$. The timer stays public, because the sliced policy schedules by it.

\paragraph{Output} A pool output is a six-tuple of options, one per slot, and
\m{out_equiv} is the corresponding nest of \m{eqpair}s (\cref{tab:out-equiv}). Its two
columns come apart only for the two secret handlers, and that is where scheduling
secrecy lives. The payload equivalence hides a slot's \emph{value}. Whether the slot
\emph{ran} is decided by the \code{eqmaybe} instance around it, which fixes who can see
\code{None}. Under \m{eqmaybe_all} the attacker tells \code{Some} from \code{None}, and so
sees that the private user process ran; under \m{eqmaybe_nonbot} a handler run and
a step in which nothing happened are the same observation.

\begin{table}[t]
\centering
\caption{\texttt{out\_equiv}, one component per slot: a wrapper from the
  \texttt{eqmaybe} family around a public or private payload equivalence. The last
  two columns say what the attacker at $\Lo$ sees.}
\label{tab:out-equiv}
\footnotesize
\begin{tabular}{@{}lllll@{}}
\toprule
slot & wrapper & payload & value & ran? \\
\midrule
public user & \m{eqmaybe_all} & public & visible & visible \\
private user & \m{eqmaybe_all} & private & hidden & visible \\
scheduler & \m{eqmaybe_all} & public & visible & visible \\
NOP & \m{eqmaybe_nonbot} & private & hidden & hidden \\
disk & \m{eqmaybe_nonbot} & private & hidden & hidden \\
timer & \m{eqmaybe_all} & public & visible & visible \\
\bottomrule
\end{tabular}
\end{table}

The NOP slot must be classified like the disk slot: a slice run is filled by one or
the other, so the two must look alike.

\subsection{Compositional structure}

The proof follows the term of \m{model_sliced} from the outside in, applying one
composition theorem per constructor (\cref{thm:composition}): \m{map_NI} for the outer
rewiring, \m{loop_NI} for the feedback, \m{sta_NI} for the global state, \m{maybe_NI} for
the idling pool, \m{par_NI} to split the pool into slots, \m{swi_NI} for each slot's
gate, and at the leaves \m{out_NI} for the handler processes and the hypotheses of
\cref{thm:main} for the scheduler and user processes. Since each theorem derives the
composite's equivalences from those of its parts, the shape of \m{out_equiv} is forced:
\m{par_NI} imposes its nest of \m{eqpair}s and \m{swi_NI} its \code{eqmaybe} instances.

Two of these steps carry substantial side conditions: the gate obligations of
\m{swi_NI} (\cref{sec:gating}) and the state obligations of \m{sta_NI}
(\cref{sec:state-equiv,sec:fvni}).

\subsection{Gating a slot}\label{sec:gating}

Each slot is a \code{swi} whose gate bit is computed from the pid by a \code{map}. Two
obligations remain per slot, both parameterised by an equivalence $\eqv{B}$ on the gate
bit that the proof may choose. \m{map_NI} requires the gate map to be unobservable-preserving,
sending unobservable pids to unobservable bits (\m{f_PU}), and \m{swi_NI} requires, at every level, that the
observer be aware of the gate or oblivious to the slot:
\[
\ell \in \Aw(\mathsf{true}, \eqv{B}) \;\cup\;
  \Ob(\mathsf{swi}\;b\;p,\; \mathsf{eqmaybe\_aware}(\eqv{O},\eqv{B})).
\]
Awareness means a bit related to \code{true} is \code{true} and observable, so the observer can
trust the gate. Obliviousness means there is a reference output $o_0$ such that every
output of the slot is related to it: whatever the input does, the observer sees the
same class. This is weaker than the original condition (\cref{sec:departures}), and the
argument below needs the weaker form.

At $\Lo$ the two pids of the secret handlers are unobservable. For a slot neither of
them selects, both map to \code{false}, so a $\eqv{B}$ in which \code{false} alone is unobservable
satisfies \m{f_PU}; we call it \m{false_equiv}. Under \m{false_equiv} the observer is aware
at every level, and the obligation is discharged without a case split. This covers
the public and private user slots, the scheduler and the timer handler.

For the disk and NOP slots, one secret pid maps to \code{true} and the other to \code{false},
so \m{f_PU} demands that both bits be unobservable, which leaves only \m{private_equiv}.
But then \code{true} is unobservable at $\Lo$ and awareness fails there. At $\Hi$
awareness still holds; at $\Lo$ the proof uses obliviousness, with \code{None} as the
reference output. A step where the slot is gated off emits \code{None}, and a step where
it runs emits \code{Some o}, so the obligation is that \code{None} and \code{Some o} are related at
$\Lo$. Under \m{eqmaybe_aware} that holds exactly when the gate is hidden and the
payload unobservable, which the \m{private_equiv} payload of \cref{tab:out-equiv}
provides. The line therefore falls between the two secret handlers and everything
else; the private user process sits on the public side, its secrecy carried by its
payload and never by its gate.

\subsection{The state equivalence}\label{sec:state-equiv}

\m{sta_NI} imposes two conditions on the state update, and they pull in opposite
directions.
\begin{itemize}
\item \m{f_EP} requires that an unobservable input not visibly move the state. A disk
  interrupt is unobservable at $\Lo$ and sets the disk pending bit, so that bit
  must be private. (The event equivalence \code{eqsum_L in_equiv out_equiv} makes only
  interrupts unobservable, so \m{f_EP} constrains \m{set_pending} alone.)
\item \m{fv_NI} requires that related states be mapped to related states.
  \m{initiate_next} branches on the mask bits, so two related states that disagreed
  there would take different branches; the mask bits must be public.
\end{itemize}
Both hold only if the interrupt controller is classified per bit: the disk and NOP
pending bits are private and every mask bit is public (\cref{tab:state-equiv}).
The same constraint explains why \m{set_pending} sets one bit and stops: any decision
taken on an input step, such as starting a handler or moving the slice counter,
would disturb a public field. Everything an interrupt causes is deferred to a later
output step, where the input is no longer the thing being varied.

\begin{table}[t]
\centering
\caption{\texttt{state\_type\_equiv}, field by field.}
\label{tab:state-equiv}
\footnotesize
\begin{tabularx}{\columnwidth}{@{}>{\raggedright\arraybackslash}p{0.25\columnwidth}L>{\raggedright\arraybackslash}p{0.24\columnwidth}@{}}
\toprule
field & equivalence & at $\Lo$ \\
\midrule
current pid & \code{eqsum private_equiv public_equiv} & which secret handler is hidden \\
previous pid, \m{re_sched}, \m{ir_count} & \m{public_equiv} & visible \\
NOP and disk bits & \m{private_pending_equiv} & pending hidden, mask visible \\
timer bits & \m{public_bits_equiv} & visible \\
\bottomrule
\end{tabularx}
\end{table}

The pid type is $\mathsf{Sum}\;\mathsf{Bool}\;\mathsf{Nat}$. Its \code{inl} side
holds exactly the two secret handlers, the disk handler and the NOP handler, and its
\code{inr} side everything else, the timer handler included. Drawing the line there
lets the proof classify a whole pid in one place.

The current pid is where the policy does its work. Since \m{eqsum} never relates an
\code{inl} to an \code{inr}, the fact \emph{that} a handler is running stays visible, as does
every public pid. What \m{private_equiv} hides is the bit inside \code{inl}: at $\Lo$ the
disk handler and the NOP handler look alike. Across two $\Lo$-related states, that
bit and the two private pending bits are the only things allowed to differ.

\subsection{Closure under the transition}\label{sec:fvni}

The core obligation is that \m{state_type_equiv} is closed under the whole transition.
The transition is a sum, \m{set_pending} on interrupts and a pipeline of stages on
outputs, and the event equivalence never relates an interrupt to an output. So the
goal splits along the sum, and a composition lemma (\m{fv_NI_comp}) breaks the
pipeline into one goal per stage:
\begin{itemize}
\item \m{set_pending}: related interrupts set related pending bits;
\item \m{check_scheduler}: the scheduler's pid is public, so the update agrees;
\item \code{service_policy}: the slice logic reads only the public counter and
  masks, and writes the same values into both states;
\item \m{initiate_next}: related states must make the same choice of what runs
  next.
\end{itemize}
Each goal is self-contained, which is the point of splitting. The cost is that a
stage must be proved for \emph{all} pairs of related states, since it cannot assume
its input came from the stage before. That loses the invariant of the time slice,
and the last stage cannot be proved without it.

\subsection{Restoring the invariant}\label{sec:restore}

Two facts hold of every state \m{model_sliced} reaches. While a slice is live and no
handler is running, the NOP handler is pending, the disk and NOP masks are clear and
the timer is masked; once the slice is over, the disk and NOP handlers are masked.
And the disk and NOP masks always agree. The sliced policy maintains both
(\cref{sec:sliced}).

The function \m{restore_invariant} writes these facts into the state, just before
\m{initiate_next} chooses. It reads only the public counter, so it sets the same bits in
two related states, and afterwards the deciding mask bits agree and related states
take the same branch. This is what discharges the \m{initiate_next} goal.

We prove that it never changes a state the model reaches.

\begin{theorem}[\m{restore_invariant_noop}]\label{thm:noop}
Let $\mathit{runtime} \neq 1$. For every sequence of events $es$ from the initial
state of \m{model_sliced}, and every output $o$, let $w$ be the state that
\m{initiate_next} receives on $o$ after $es$. If not all masks are set in $w$, then
$\mathsf{restore\_invariant}\;w = w$.
\end{theorem}

The events $es$ are arbitrary, so the theorem holds whatever the pool emits. The proof states the invariant
between steps, shows that every event preserves it, and shows that at the choice
point a slightly weaker invariant holds, on which \m{restore_invariant} is the
identity. One small lemma per stage of the pipeline does the bookkeeping.

The hypothesis $\mathit{runtime} \neq 1$ exposes a quirk of the model. The boundary
test reads the counter before the stage that ticks it, while \m{restore_invariant}
reads it after. With one-step handlers the last boundary inside a slice ticks the
counter to zero in the same step, so the slice already reads as over, and
\m{restore_invariant} masks the handler that boundary had just enabled. Non-interference
is unaffected (\cref{thm:main} has no such hypothesis), but at $\mathit{runtime} = 1$
every slice loses its last handler run.

An earlier version of the model relied on \m{restore_invariant} for behaviour: it was
the only thing that queued the NOP filler and the only thing that returned control
to the scheduler after a slice. Proving \cref{thm:noop} required moving both into
\code{service_policy}, where they now appear as part of the policy
(\cref{sec:sliced}). The proof of \cref{thm:main} needed only two small lemmas to
follow the change.
